\documentclass[12pt,a4paper]{report}

\usepackage[
  top=1in,
  bottom=1in,
  left=1.2in,
  right=1in
]{geometry}
\usepackage[T1]{fontenc}
\usepackage{lmodern}
\usepackage{graphicx}
\usepackage{float}
\usepackage{xcolor}
\usepackage{booktabs}
\usepackage{tabularx}
\usepackage{array}
\usepackage{enumitem}
\usepackage{setspace}
\usepackage{listings}
\usepackage{caption}
\usepackage[hidelinks]{hyperref}
\usepackage{fancyhdr}

\graphicspath{{screenshots/}}
\setlength{\parskip}{0.45em}
\setlength{\parindent}{0pt}
\setlist[itemize]{leftmargin=1.5em,itemsep=2pt}
\setlist[enumerate]{leftmargin=1.8em,itemsep=2pt}
\renewcommand{\arraystretch}{1.2}
\newcolumntype{Y}{>{\raggedright\arraybackslash}X}
\newcommand{\INR}{INR\,}
\newcommand{\screenshot}[3]{%
  \begin{figure}[H]
    \centering
    \includegraphics[width=\linewidth,height=#3\textheight,keepaspectratio]{#1}
    \caption{#2}
    \label{fig:#1}
  \end{figure}
}

\lstset{
  basicstyle=\ttfamily\small,
  frame=single,
  breaklines=true,
  columns=fullflexible,
  keepspaces=true,
  showstringspaces=false
}

\pagestyle{fancy}
\fancyhf{}
\fancyhead[L]{Cairn: Governed Agentic Commerce}
\fancyhead[R]{Project Report}
\fancyfoot[C]{\thepage}
\setlength{\headheight}{15pt}

\begin{document}

\begin{titlepage}
  \centering
  \vspace*{0.8cm}
  {\Large\bfseries RV UNIVERSITY\par}
  \vspace{0.2cm}
  {\small RV Vidyanikethan Post, 8th Mile, Mysuru Road, Bengaluru -- 560 059\par}
  \vspace{0.65cm}
  {\large\bfseries SCHOOL OF COMPUTER SCIENCE AND ENGINEERING\par}
  \vspace{1.3cm}
  {\Large Project Report\par}
  \vspace{0.5cm}
  {\LARGE\bfseries Cairn: A Governed Agentic-Commerce Platform\par}
  \vspace{0.3cm}
  {\large Buyer and Merchant Workflows with a Shared Deterministic Governance Layer\par}
  \vspace{1.0cm}
  {\large CS4951 -- Major Project\par}
  \vspace{0.25cm}
  {\large Bachelor of Science (Honours)\par}
  \vspace{1.0cm}
  {\large\bfseries Submitted by\par}
  \vspace{0.3cm}
  \begin{tabular}{ll}
    Hampana SS & 1RVU23BSC050 \\
    Shamitha Reddy N & 1RVU23BSC131 \\
    Mohnish G Naidu & 1RVU23BSC084 \\
    Kotta Sivaganesh & 1RVU23BSC069
  \end{tabular}
  \vspace{0.9cm}

  {\large\bfseries Under the guidance of\par}
  \vspace{0.2cm}
  Prof.\ Anoop A\\
  Assistant Professor, School of Computer Science and Engineering

  \vfill
  {\large August -- December 2026\par}
\end{titlepage}

\pagenumbering{roman}

\chapter*{Declaration}
\addcontentsline{toc}{chapter}{Declaration}
We declare that this report describes the Cairn major project undertaken by
the student team listed on the title page. The implementation status and
limitations are reported as observed in the repository at the time of writing.
Any institution-specific declaration wording and signatures should be reviewed
and completed by the authors and project guide before formal submission.

\chapter*{Acknowledgement}
\addcontentsline{toc}{chapter}{Acknowledgement}
We thank our project guide, Prof.\ Anoop A, and the School of Computer Science
and Engineering at RV University for their guidance and support. We also thank
our families and peers for their encouragement during the project.

\chapter*{Abstract}
\addcontentsline{toc}{chapter}{Abstract}
Cairn is a buyer--merchant e-commerce prototype designed around the principle
that AI agents may interpret requests and propose actions, while explicit
software rules and human approvals control consequential commerce operations.
The project combines a React and Vite interface with a FastAPI backend,
SQLite-backed marketplace data, vendor discovery and allocation, order
life-cycle operations, scoped notifications, and an optional local Ollama
integration. The buyer experience includes governed shopping concepts,
spending preferences, order status and notifications; the merchant experience
provides an order inbox and operational views.

The backend prototype includes ten synthetic Bengaluru vendors, geolocation
helpers, ranked inventory allocation and reservation, merchant order actions,
buyer/vendor notifications, server-sent events (SSE), SLA records, a sandbox
payment adapter and verifiable receipt support. The automated backend suite
contains 37 passing tests and the frontend production build completes. These
results establish a functional local prototype, not production readiness.
Important remaining work includes authenticated identities and ownership
checks, connecting allocation to governance decisions, transactional and
idempotent order writes, validated end-to-end location data, durable
notifications, and production database migrations.

\textbf{Keywords:} agentic commerce, deterministic governance, human approval,
marketplace allocation, explainability, order life cycle.

\tableofcontents
\listoffigures
\listoftables
\clearpage
\pagenumbering{arabic}

\chapter{Introduction}

\section{Background}
Conversational AI can make product discovery easier, but an agent that is
allowed to initiate purchases also creates risks: spending beyond a buyer's
intent, acting on untrusted merchant data, or making decisions that are hard to
explain. Cairn explores a design in which an agent can help with discovery
without being treated as the authority for payment, approval, or inventory
decisions.

\section{Problem statement}
The project addresses the need for a two-sided commerce workflow in which
buyer intent, seller fulfillment, and marketplace actions are visible and
traceable. A useful prototype must demonstrate both the buyer journey and the
merchant order journey, connect them to backend data, and disclose when a
workflow still relies on demo data or has unresolved security boundaries.

\section{Objectives}
\begin{enumerate}
  \item Provide buyer and merchant interfaces with consistent navigation and
        accessible interaction patterns.
  \item Keep purchase authorization and policy checks in deterministic
        application logic rather than delegating them to a language model.
  \item Model vendors, inventory, delivery areas, reservations, and order
        allocation for a local marketplace prototype.
  \item Support merchant order actions and reflect their updates in buyer
        orders and notifications.
  \item Add focused backend tests and document the current implementation and
        its production blockers.
\end{enumerate}

\section{Scope and status}
This report covers the local prototype in the repository. The ten seeded
Bengaluru vendors are synthetic demonstration records. The payment integration
is a sandbox adapter, and Ollama is an optional local inference service. The
prototype must not be used to process real customer payments or represent the
seeded vendors as real businesses.

\chapter{System Design}

\section{Design principles}
\begin{itemize}
  \item \textbf{Agents propose; application rules decide.} Language models can
        help interpret or summarize requests, but purchase and allocation
        decisions require explicit deterministic checks.
  \item \textbf{Human control.} Buyer approval and opt-in notification sound
        provide user control where the interface supports them.
  \item \textbf{Explainability.} Order status, governance concepts and
        receipt information should be presented in user-understandable terms.
  \item \textbf{Fail honestly.} Missing backend data or unavailable services
        should not be presented as verified outcomes.
  \item \textbf{Two-sided synchronization.} Merchant order actions should
        update the corresponding buyer view through shared backend state.
\end{itemize}

\section{Architecture}
The frontend is a React application served by Vite. It calls a FastAPI
backend, which exposes governance, marketplace, order, notification, payment
sandbox and telemetry operations. SQLite stores local prototype data. The
backend contains separate modules for vendor storage, geolocation, allocation,
order life cycle, notifications, payment simulation and receipts. The LLM
client can call a local Ollama server; deterministic fallback behavior is
available when that service cannot be reached.

\begin{table}[H]
  \centering
  \caption{Main implementation components}
  \label{tab:components}
  \begin{tabularx}{\textwidth}{p{0.23\textwidth}YY}
    \toprule
    \textbf{Layer} & \textbf{Technology or module} & \textbf{Responsibility} \\
    \midrule
    User interface & React, Vite, CSS & Buyer and merchant screens,
      navigation, order status and notifications \\
    API & FastAPI, Uvicorn & Health, governance, vendor, allocation,
      order and notification endpoints \\
    Governance & Python modules & Authorization, policy checks, conflict
      detection, risk scoring and explanations \\
    Marketplace & \texttt{backend/marketplace/} & Vendor data, location
      helpers, allocation, reservations, notifications and order service \\
    Persistence & SQLite & Local prototype records and test databases \\
    Language model & Ollama client & Optional local intent/agent assistance
      with a deterministic fallback \\
    \bottomrule
  \end{tabularx}
\end{table}

\section{Request and order flow}
\begin{enumerate}
  \item The buyer describes a need or opens a marketplace workflow.
  \item The application obtains vendor and inventory candidates using the
        available location and product information.
  \item The allocation service filters and ranks candidates and attempts to
        reserve stock.
  \item The merchant reviews an allocated order and may accept, prepare,
        reject, propose a price change, or provide shipment details.
  \item The backend updates order data and creates buyer/vendor notifications.
        The interfaces refresh through scoped SSE and polling fallback.
  \item The buyer reviews the updated order state. Notification sound is
        opt-in and requires the browser tab to remain open.
\end{enumerate}

\textbf{Important boundary:} the current marketplace allocation endpoint is
not yet routed through the deterministic governance engine. Therefore the
prototype's presence of both governance modules and allocation APIs does not
mean that every allocation is currently governed end-to-end.

\chapter{Implementation}

\section{Buyer experience}
The buyer side includes role selection, sign-in options, Explore with Aria,
spending-rule controls, order status, and a notification centre. The interface
is designed around a calm visual style and explicit controls. The eight
screenshots included with this report are UI captures; screenshots alone do
not verify that an LLM response or purchase took place.

\section{Merchant experience}
The merchant side includes dashboard and store-health views, a vendor order
inbox, and actions for accepting, preparing, rejecting, reallocating and
dispatching an order. Shipment confirmation requires a carrier and tracking
reference. Lifecycle changes are shared with the buyer order view in the
prototype.

\section{Vendor allocation and geolocation}
The marketplace store seeds ten synthetic vendors in Bengaluru when the
vendor table is empty. The geolocation module provides haversine distance and
service-radius helpers. Allocation filters and ranks possible vendor
inventory, then reserves stock. Known gaps remain: the buyer flow often sends
locality rather than validated coordinates, custom location entry is not
geospatially resolved correctly, and price bounds and vendor-open status are
not consistently enforced as eligibility conditions.

\section{Notifications and service levels}
Buyer/vendor notifications are persisted and scoped SSE updates are supported,
with periodic polling as a refresh fallback. The event delivery mechanism is
process-local, has no replay or cross-worker delivery, and is not an
authenticated push service. SLA records are created for order acknowledgement
and dispatch, but breach checks are manually triggered rather than managed by
a supervised scheduler.

\section{Ollama and API key requirements}
The local Ollama path does not require a provider API key. A developer can
install Ollama, pull the configured model, start the local service, and run
the backend with the repository's settings. The README documents an example
using \texttt{llama3.2:1b}. An external hosted LLM provider would require its
own credential, but no such key should be committed to the repository.
Regardless of model availability, authorization and commerce policy should
remain deterministic.

\section{Sandbox payments and receipts}
The payment module provides a deterministic sandbox adapter for create,
confirm and refund operations, and the receipt module supports HMAC-SHA256
signing and verification. These are prototype facilities only: no real
payment gateway, webhook verification, bank settlement, or refund policy is
connected.

\chapter{Screenshots and User Interface}
The screenshots below are stored in the repository's \texttt{screenshots/}
folder and can be uploaded with the report to Overleaf. The three merchant
captures were taken while backend data was unavailable; they document the UI
state, not live order, inventory, SLA or governance measurements.

\screenshot{01-role-selection.png}
  {Role selection landing page for buyer and merchant experiences}{0.70}

\screenshot{02-buyer-explore.png}
  {Buyer Explore screen with the Aria conversational entry point}{0.70}

\screenshot{03-buyer-sign-in.png}
  {Buyer sign-in screen with Google, demo and guest options}{0.70}

\screenshot{04-aria-governance-processing.png}
  {Aria request and governance processing UI state; not evidence of a completed order}{0.70}

\screenshot{05-merchant-buyer-orders.png}
  {Merchant buyer-orders view captured while backend data was unavailable}{0.70}

\screenshot{06-merchant-dashboard.png}
  {Merchant dashboard layout captured while backend data was unavailable}{0.70}

\screenshot{07-merchant-store-health-audit.png}
  {Store health and audit layout captured while backend data was unavailable}{0.70}

\screenshot{08-buyer-spending-rules.png}
  {Buyer spending rules and safeguards interface}{0.70}

\chapter{Testing and Results}

\section{Validation performed}
The repository validation recorded for this prototype is summarized in
Table~\ref{tab:validation}. The backend tests use isolated SQLite test
databases. The results are evidence for the tested prototype paths; they do
not establish production security, concurrent inventory correctness, or
complete end-to-end reliability.

\begin{table}[H]
  \centering
  \caption{Prototype validation results}
  \label{tab:validation}
  \begin{tabularx}{\textwidth}{p{0.30\textwidth}p{0.19\textwidth}Y}
    \toprule
    \textbf{Check} & \textbf{Result} & \textbf{Interpretation} \\
    \midrule
    Backend test suite & 37 passed & Covers API, governance, LLM client and
      marketplace lifecycle paths in isolated test databases \\
    Frontend production build & Passed & Vite generated a production build \\
    Frontend lint & Passed with warnings & Existing lint warnings remain and
      should be addressed separately \\
    \bottomrule
  \end{tabularx}
\end{table}

\section{Tested marketplace behaviors}
Focused tests cover vendor allocation, notifications, accept/prepare/dispatch,
price proposals, and reject/reallocate/accept. These tests do not prove
correctness under concurrent allocation load or failures across multiple
database writes. Browser end-to-end tests for the full buyer-to-merchant
workflow remain future work.

\section{Limitations and production blockers}
\begin{enumerate}
  \item \textbf{Identity and authorization:} buyer and vendor identity is
        demo-grade. Caller-supplied identifiers are not a substitute for
        authenticated principals, and order ownership is not enforced
        consistently.
  \item \textbf{Governance integration:} allocation can create an order
        without invoking the deterministic governance engine or approval
        policy.
  \item \textbf{Atomicity and idempotency:} reservation, order, SLA and
        notification writes are separate operations. Retries and partial
        failures can produce inconsistent records or stranded reservations.
  \item \textbf{Location:} buyer coordinates are not consistently used and
        validated across discovery and allocation.
  \item \textbf{Notifications and SLAs:} SSE is process-local and unauthenticated;
        there is no replay, durable cross-worker bus, scheduled SLA worker or
        closed-app push delivery.
  \item \textbf{Database portability:} SQLite-specific SQL is present.
        PostgreSQL compatibility and migration behavior have not been
        established.
  \item \textbf{External services:} there are no real payment, email, SMS,
        WhatsApp or push-notification integrations. Seed vendors and data are
        synthetic.
\end{enumerate}

\chapter{Conclusion and Future Work}

\section{Conclusion}
Cairn demonstrates a two-sided agentic-commerce prototype with buyer and
merchant workflows, backend marketplace data, vendor allocation, order
life-cycle actions, synchronized in-app notifications, optional local Ollama
assistance, and governance-oriented explanations. The implementation has
focused automated backend coverage and builds successfully. The current
result is suitable for local demonstration and further development, but not
for real customer orders, real payments, or production vendor onboarding.

\section{Recommended next steps}
\begin{enumerate}
  \item Add authenticated buyer and vendor principals and enforce ownership
        checks for all order and notification operations.
  \item Route allocation and purchase actions through the deterministic
        governance and approval workflow.
  \item Make reservation, allocation, order, SLA and notification writes
        transactional and idempotent; test retries and failure recovery.
  \item Validate location input and use the same resolved coordinates for
        discovery and allocation.
  \item Add authenticated durable notification delivery and a scheduled SLA
        processor.
  \item Introduce versioned migrations and test the intended production
        database dialect.
  \item Add browser end-to-end and concurrent-allocation tests before
        integrating real vendors or payment services.
\end{enumerate}

\chapter*{References}
\addcontentsline{toc}{chapter}{References}
\begin{enumerate}
  \item React, ``React Documentation.'' Available:
        \url{https://react.dev/}
  \item FastAPI, ``FastAPI Documentation.'' Available:
        \url{https://fastapi.tiangolo.com/}
  \item Ollama, ``Ollama Documentation.'' Available:
        \url{https://docs.ollama.com/}
  \item Amazon, ``Seller Central.'' Available:
        \url{https://sell.amazon.com/tools/seller-central}
  \item Cairn project repository. Available:
        \url{https://github.com/Mohnish-140605/agentic_governance_ecommerce}
\end{enumerate}

\end{document}
